%!TEX root = ../hutbthesis_main.tex

\chapter{绪论}

\section{研究背景与意义}

海洋资源开发、水下工程作业、水下探测与搜救、水下装备操作训练等场景越来越多地依赖水下机器人与水下传感系统。与陆地环境相比，水下工作环境具有噪声复杂、信号传播延迟高、载体易受水流影响产生漂移、光线遮挡严重、水体能见度低等特点，光学相机、激光、雷达等常用感知手段在水下的作用距离与成像质量明显下降，声呐成为水下远距离感知的主要手段，但声呐成像存在分辨率有限、多径混响干扰、目标轮廓模糊等问题\cite{ref1}。与此同时，水下信息传输主要依靠水声通信，水声信道带宽窄、多径效应强、传播时延大、易受丢包与多普勒漂移影响，使得传感器数据回传与遥控指令链路的稳定性难以保证\cite{ref2}。

在真实水下环境中开展感知算法与通信策略研究，存在试验成本高、风险大、水文条件不可重复、样本采集困难等问题，难以支撑算法的反复调试与对比验证。借助仿真平台在虚拟环境中复现水下噪声、延迟、漂移、遮挡、低能见度与声呐成像等典型因素，并模拟传感器数据与通信链路的传输过程，可以在安全、低成本、可重复的条件下完成算法验证与操作人员训练，具有明确的工程应用价值与理论意义\cite{ref3}。

本课题面向水下工作环境，研究传感器与通信链路的仿真方法，构建一套可配置环境参数、可复现典型水下干扰、支持算法验证与操作训练的仿真系统。【待补充：结合导师给定的具体应用背景，进一步说明本课题对应的实际需求与研究价值。】

\section{国内外研究现状}

\subsection{水下环境与传感器仿真研究现状}

【待补充：梳理水下环境建模、光学/水声传感器仿真、声呐成像模拟等方向的国内外代表性工作，正文引用须与 references.tex 中的条目一一对应。】

在水下感知仿真方面，现有研究主要围绕水体光学衰减、悬浮颗粒与浑浊度建模、声呐波束与混响仿真、传感器噪声与误差建模等方向展开\cite{ref1}。【待补充：国内研究进展，列出代表性团队、方法与结论。】

\subsection{水声通信链路仿真研究现状}

【待补充：梳理水声信道建模、时延与多径模拟、带宽与丢包仿真、通信协议与链路质量评估等方向的研究现状\cite{ref2}。】

\subsection{研究述评}

综合现有工作可以看出，水下传感器仿真与水声通信链路仿真已分别形成若干方法，但面向操作训练与算法验证的一体化仿真平台仍存在不足：环境干扰因素（噪声、延迟、漂移、遮挡、低能见度、声呐成像）的配置不够直观，传感器仿真与通信链路仿真相互割裂，缺少统一的测试与评价手段。【待补充：结合文献进一步归纳现有研究的局限，明确本课题的切入点。】

\section{研究内容与创新}

本课题围绕“基于水下工作环境的传感器与通信链路仿真系统”展开，主要研究内容如下：

1. \textbf{水下工作环境建模}：构建可配置的水下场景，对噪声、延迟、漂移、遮挡、低能见度等环境因素进行参数化建模。
2. \textbf{传感器感知仿真}：模拟光学/水声等传感器在上述干扰条件下的感知结果，包括声呐成像结果的生成与显示。
3. \textbf{通信链路仿真}：模拟水声通信链路的延迟、带宽、丢包、多径等特性，对传感器数据与控制指令的传输过程进行建模。
4. \textbf{系统交互与测试评估}：提供面向算法验证和操作训练的交互界面，在指定仿真场景下完成功能测试与性能测试，并记录可对比的实验结果。

本课题的创新点【待系统实现后据实提炼，未测试通过的功能不得写入创新点】：

1. 【待补充：环境建模或传感器仿真方面的改进点。】
2. 【待补充：通信链路仿真或多因素耦合方面的改进点。】
3. 【待补充：面向操作训练的系统设计或评价方法方面的特点。】

\section{论文组织架构}

本论文按照“需求分析—环境与传感器建模—通信链路建模—系统设计与实现—测试与分析—总结展望”的顺序展开，各章节安排如下：

1. \textbf{绪论}：介绍研究背景与意义、国内外研究现状、研究内容与论文组织结构。
2. \textbf{相关技术与理论基础}：【待补充：仿真平台、水下传感器、水声通信、声呐成像等相关技术概述。】
3. \textbf{系统需求分析与总体设计}：【待补充：需求分析、系统架构、模块划分与数据流设计。】
4. \textbf{环境与传感器、通信链路建模实现}：【待补充：各仿真模块的具体实现。】
5. \textbf{系统实现与交互界面}：【待补充：系统编码实现、入口程序与交互界面。】
6. \textbf{测试与结果分析}：【待补充：功能测试、性能测试、典型案例与失败案例分析。】
7. \textbf{总结与展望}：【待补充：工作总结、存在不足与后续改进方向。】
